\section{对角线引理：把代入做成自指}\label{sec:diagonal}

\begin{definition}[自代入函数]\label{def:self-substitution}
按课堂板书，先将代入写成三元函数 \(\Subst(a,\goedel{x},t)\)，再定义
\[
S(a,b)=\Subst(a,\goedel{x},\Num(b)),\qquad s(u)=S(u,u).
\]
若 \(u=\goedel{\theta(x)}\)，则 \(s(u)=\goedel{\theta(\num u)}\)。函数 \(s\) 原始递归，故 \(\PA\) 可表示并逐个数词计算它。
\end{definition}

\begin{theorem}[对角线引理]\label{thm:diagonal}
对每个只含自由变量 \(x\) 的公式 \(\varphi(x)\)，存在句子 \(\psi\) 使
\[
\PA\vdash \psi\leftrightarrow \varphi(\num{\goedel{\psi}}).
\]
\end{theorem}

\begin{proof}
令
\[
\theta(x)\equiv \varphi(s(x)),\qquad
m=\goedel{\theta(x)},\qquad
\psi\equiv\theta(\num m).
\]
由 \(s\) 的定义，
\[
s(m)=\goedel{\theta(\num m)}=\goedel{\psi}.
\]
因 \(s\) 可表示，\(\PA\) 证明 \(s(\num m)=\num{\goedel{\psi}}\)。于是
\[
\PA\vdash \psi\leftrightarrow \varphi(\num{\goedel{\psi}}).
\]
\end{proof}

\begin{remark}[与课堂板书写法的对应]\label{rem:blackboard-diagonal}
课堂板书写的是
\[
\theta(x)\equiv\varphi(S(x,x)),\quad
m\equiv\goedel{\theta(x)},\quad
\psi\equiv\theta(\num m)
=\varphi(S(\goedel{\theta(x)},\goedel{\theta(x)})).
\]
关键不在代数技巧，而在层次区分：\(\goedel{\theta(x)}\) 是元语言中的数，\(\num m\) 是对象语言中的数词，\(\theta(\num m)\) 又是公式。代入函数作用在编码上，而不是作用在直观文本上。
\end{remark}
